\section{Conclusions}
\label{sec:conclusion}

Public practitioner discussions contain enough operational detail to support a rigorous qualitative pilot and bounded quantitative program for connected laboratories. Their strongest scientific use is not estimation of field-wide failure rates. It is the identification of measurable engineering objects, missing contracts, evidence boundaries, and prospective outcomes.

The study identifies ten bottlenecks across ecosystem knowledge, interfaces and representations, runtime coordination, evaluation, and AI-specific context. Their common structure concerns transitions from one representation to an operational commitment. Knowledge must become a maintained artifact. Device access must become a testable interface. Source must become an identified deployment. Resource data must become a validated physical model. Observations must become linked evidence. Interrupted execution must become a verified recovery decision. Generated methods must become experimentally evaluated procedures.

The strongest next empirical study concerns observability and recovery. The pilot supplies the strongest descriptive association and direct incident mechanisms. A prospective incident dataset can test whether command--state--material linkage predicts time to a verified safe state, material salvage, final disposition completeness, and recurrence. The released troubleshooting schema, knowledge index, interface registry design, and event schema provide immediate community infrastructure and the measurement substrate for that work.
